\chapter{Introduction}
\label{ch:introduction}

\begin{chapterabstract}
This chapter situates the thesis within the rapidly evolving landscape of
embodied artificial intelligence, where vision-language models are
increasingly deployed on resource-constrained robotic platforms. It
identifies the central inefficiency that motivates this work --- the
repeated invocation of expensive multimodal reasoning even when the
semantic content of a scene is essentially unchanged --- and frames the
research question, objectives, and contributions around the reduction of
this \emph{semantic reasoning redundancy}.
\end{chapterabstract}

\section{Background and Contemporary Context}
\label{sec:intro-background}

The convergence of large-scale vision-language models (VLMs) with mobile
robotic platforms has opened a new frontier in embodied perception. Models
such as CLIP~\cite{Radford2021CLIP}, BLIP-2~\cite{Li2023BLIP2}, and
LLaVA~\cite{Liu2023LLaVA} couple visual encoders with language reasoning to
produce open-vocabulary scene descriptions, answer questions about an
environment, and ground natural-language instructions in physical
affordances. When integrated into robotic control stacks --- as in
RT-2~\cite{Brohan2023RT2}, SayCan~\cite{Ahn2022SayCan}, and
OpenVLA~\cite{Kim2024OpenVLA} --- these models allow a robot to interpret its
surroundings with a flexibility that hand-engineered perception pipelines
cannot match.

This expressive power, however, carries a substantial computational price.
A single forward pass through a modern VLM may involve billions of
floating-point operations, hundreds of milliseconds of latency, and several
gigabytes of accelerator memory. On the consumer-grade and embedded
accelerators typical of edge robots --- as opposed to the data-centre GPUs on
which these models are trained --- such costs translate directly into reduced
battery life, thermal throttling, and an inability to maintain real-time
perception rates. The tension between the semantic richness of VLMs and the
strict energy, latency, and memory envelopes of edge robotics is the
defining systems challenge addressed by this thesis.

Crucially, a robot perceiving a continuous video stream does not encounter a
sequence of independent, equally informative images. Consecutive frames are
strongly correlated: a robot pausing in a corridor, waiting for a door to
open, or observing a static workspace produces long runs of frames whose
\emph{semantic} content barely changes even as raw pixels fluctuate with
sensor noise, illumination, and minor camera motion. The opportunity that
this thesis exploits is the gap between \emph{pixel-level} change and
\emph{semantic-level} change: the former is frequent and cheap to measure,
whereas the latter is rare and is what actually warrants fresh multimodal
reasoning.

\section{The Core Problem: Semantic Reasoning Redundancy}
\label{sec:intro-problem}

Contemporary robotic VLM systems overwhelmingly adopt a \emph{stateless,
per-frame} invocation discipline: each incoming frame, or every $K$-th frame
under a fixed sampling schedule, is passed to the VLM regardless of whether
the scene has meaningfully evolved since the last invocation. This design is
simple and reproducible, but it is profoundly wasteful. When a robot stares
at an unchanging scene for several seconds, a per-frame policy re-derives the
same caption or answer dozens of times, consuming accelerator cycles and
energy to produce information the system already possesses.

We term this inefficiency \emph{semantic reasoning redundancy}: the repeated
expenditure of heavyweight multimodal compute on inputs whose semantic state
is, for all practical purposes, identical to a previously reasoned-over
state. It is conceptually distinct from the redundancies that prior efficiency
research targets. Pixel- and frame-level redundancy --- addressed by adaptive
frame sampling~\cite{Wu2019AdaFrame}, salient-clip
selection~\cite{Korbar2019SCSampler}, and skip
convolutions~\cite{Habibian2021Skip} --- concerns whether the raw signal has
changed. Parameter redundancy --- addressed by
quantization~\cite{Dettmers2022LLMint8}, pruning and
compression~\cite{Han2016DeepCompression}, and
distillation~\cite{Hinton2015Distillation} --- concerns whether a model is
larger than it needs to be. Neither asks the question at the heart of this
work: \emph{given what the system already understands about the scene, does
the next frame justify the cost of multimodal reasoning at all?}

Answering this question requires the system to maintain an explicit,
continuously updated notion of \emph{semantic state}, to quantify how far the
current observation has drifted from that state, and to allocate compute in
proportion to that drift. A scene that is semantically stable should be
served almost for free from memory; a scene undergoing a genuine semantic
transition --- a new object entering view, an agent leaving the frame, a task-relevant
change in configuration --- should trigger full reasoning. The absence of such a
mechanism in mainstream pipelines is the gap this thesis fills.

\section{Research Motivation and the Need for Semantic-Aware Scheduling}
\label{sec:intro-motivation}

Three observations motivate a semantic-aware scheduling approach.

\paragraph{Observation 1: Semantic change is bursty, not uniform.}
In realistic robotic deployments, semantically significant events are sparse
and clustered in time, separated by long quiescent intervals. A compute
policy that spends a fixed budget per frame is therefore mismatched to the
underlying signal: it over-provisions during quiescence and may
under-provision during bursts. A policy that tracks semantic drift can shift
budget from the former to the latter.

\paragraph{Observation 2: Cheap signals can predict the value of expensive
reasoning.}
Lightweight perception primitives that already run on edge robots --- a
nano-scale object detector~\cite{Jocher2023YOLOv8}, a multi-object
tracker~\cite{Zhang2022ByteTrack}, and a compact image
embedding~\cite{Radford2021CLIP} --- jointly summarise the scene at a small
fraction of a VLM's cost. Changes in these cheap signals are strong proxies
for the kind of semantic change that warrants invoking the VLM. The cost of
\emph{deciding} whether to reason can thus be made far smaller than the cost
of reasoning itself.

\paragraph{Observation 3: Most useful outputs are reusable.}
Because the semantic state changes slowly, the textual output of a VLM
remains valid for the duration of a stable interval. Caching the most recent
output and reusing it while drift remains below a threshold preserves the
information delivered to downstream robotic tasks while eliminating redundant
computation.

Taken together, these observations argue for a \emph{monitor-then-allocate}
architecture: a lightweight semantic monitor runs every frame, estimates a
scalar drift signal, and a scheduling policy maps that signal to a discrete
compute tier ranging from cache reuse (no VLM) to full multimodal reasoning.
This thesis designs, implements, and evaluates such an architecture.

\section{Statement of Core Research Objectives}
\label{sec:intro-objectives}

The work is organised around the following research question:

\begin{quote}
\itshape How can semantic state transitions be used to dynamically control
multimodal compute allocation in edge robotic vision systems, while
preserving the semantic fidelity of the perception delivered to downstream
tasks?
\end{quote}

This question is decomposed into four concrete objectives:

\begin{enumerate}
  \item \textbf{Formalise semantic state and semantic drift.} Define a
  composite semantic-state representation from cheap per-frame signals
  (embedding, object metadata, tracking identities) and a corresponding
  scalar drift measure that is bounded, interpretable, and sensitive to
  task-relevant change.
  \item \textbf{Design an adaptive compute-allocation policy.} Map the drift
  signal onto a small set of compute tiers via interpretable thresholds, such
  that compute is spent in proportion to semantic novelty.
  \item \textbf{Implement an edge-deployable pipeline.} Engineer a modular
  system that integrates frame filtering, detection, tracking, embedding,
  temporal memory, drift estimation, scheduling, caching, and selective VLM
  invocation, suitable for execution on commodity accelerators.
  \item \textbf{Evaluate against principled baselines.} Quantify the
  trade-off between compute saved and semantic fidelity retained, comparing
  the proposed scheduler against every-frame, uniform-sampling, motion-gating,
  and embedding-only policies using a unified metric of \emph{semantic compute
  efficiency}.
\end{enumerate}

\section{Summary of Scientific Contributions}
\label{sec:intro-contributions}

The principal contributions of this thesis are:

\begin{itemize}
  \item \textbf{A formulation of semantic reasoning redundancy} as a
  first-class efficiency target for edge VLM systems, distinct from the
  pixel-, frame-, and parameter-level redundancies addressed by prior work
  (Chapter~\ref{ch:literature}).
  \item \textbf{A composite semantic-drift model} $D_t = \alpha D_{\text{embed}}
  + \beta D_{\text{objects}} + \gamma D_{\text{track}}$ that fuses embedding
  distance, object-set novelty, and tracking-identity change into a single
  bounded signal, together with a temporal semantic memory that suppresses
  spurious drift from sensor noise (Chapter~\ref{ch:methods}).
  \item \textbf{A three-tier adaptive compute-allocation policy} that maps
  drift to cache reuse, lightweight reasoning, or full reasoning through
  interpretable thresholds $\tau_{\text{low}}$ and $\tau_{\text{high}}$
  (Chapter~\ref{ch:methods}).
  \item \textbf{A modular, edge-oriented reference implementation} and the
  \emph{Semantic Compute Efficiency} (SCE) metric that normalises retained
  semantic fidelity by compute cost, enabling fair comparison across policies
  (Chapters~\ref{ch:architecture} and~\ref{ch:experiments}).
  \item \textbf{A Semantic Observatory and a methodology for discovering
  scheduling policies.} Rather than positing a single hand-tuned heuristic,
  this thesis contributes an instrument --- the Semantic Observatory --- that
  derives a suite of semantic-evolution signals (drift, velocity,
  acceleration, volatility, scene entropy, object persistence, stability and
  transition regions) from a pipeline run, and a disciplined
  observe--hypothesise--modify-one-thing--evaluate protocol that uses these
  signals to \emph{discover and validate} scheduling policies from evidence
  (Sections~\ref{sec:meth-observatory} and~\ref{sec:arch-observatory}). The
  scientific contribution is thus a methodology for studying semantic
  evolution and deriving policies, not merely one scheduler.
\end{itemize}

\section{Thesis Document Organization}
\label{sec:intro-organization}

The remainder of this thesis is organised as follows.
Chapter~\ref{ch:literature} reviews vision-language modelling, efficient and
edge inference, temporal redundancy reduction, object detection and tracking,
and adaptive computation, positioning the present work against each.
Chapter~\ref{ch:methods} develops the proposed methodology: the semantic-state
representation, temporal memory, drift estimator, tiered scheduling policy,
and the SCE metric. Chapter~\ref{ch:architecture} describes the software
architecture and engineering design, detailing each component and the
per-frame control flow. Chapter~\ref{ch:experiments} presents the
experimental setup, baselines, evaluation metrics, quantitative results, and
an ablation study. Chapter~\ref{ch:conclusion} synthesises the findings,
discusses trade-offs and limitations, and outlines directions for future
research. Appendix~\ref{app:operational} provides a reproducibility manual.
